\section{Discussion}
By analyzing a standardized finite-element repository of 278 human pelves parameterized via an implicit continuous mapping, this study provides a population-scale quantitative evaluation of sacroiliac joint (SIJ) micromotion across locomotor and parturition-motivated mechanical regimes. The findings provide insight into how the passive architecture of the human pelvic ring accommodates the contrasting mechanical demands of bipedal weight-bearing stability and birth canal clearance \citep{Washburn1960,Ruff2017,Pavlicev2020}. Four central observations emerge: (1) SIJ kinematics display marked load-dependence, engaging self-bracing nutation under bilateral stance and pronounced asymmetric shear under single-leg stance, while parturition-motivated configurations reduce nutation; (2) obstetric mechanical proxies shift displacement from vertical craniocaudal shear toward transverse outward distraction; (3) female pelves exhibit greater translational compliance under simulated labor, an effect that is attenuated when adjusting for pelvic volume and transverse arch architecture; and (4) between-subject kinematic dispersion is governed predominantly by macro-architectural skeletal geometry, with bone mineral density heterogeneity playing a minor role ($<0.32\%$).

\subsection*{Functional morphology: form and force closure under locomotor demand}
Our findings provide quantitative support for the biomechanical framework of \emph{form and force closure} described by \citet{Vleeming2012} and \citet{Snijders1993}. Under upright bilateral stance (SP2leg), vertical ground reaction forces transmitted through both acetabula against the fixed sacrum engage physiological nutation (median $\theta_{\text{nut}} = 1.062^\circ$), rotating the sacral base anteriorly and inferiorly relative to the ilia. This sagittal tilt increases tension across the posterior interosseous and sacrotuberous ligamentous apparatus, drawing the articular surfaces into closer apposition and wedging the sacrum within the pelvic ring \citep{Snijders1993,Hammer2013,Kiapour2019}.

During single-leg stance (SP1leg), this self-bracing mechanism is subjected to asymmetric loading. Ground reaction forces transmitted through a single femoral head create an uncoupled reaction couple, driving resultant translation to $1.755$~mm and generating a substantial increase in bilateral translation asymmetry (median $2.013$~mm vs.\ $0.215$~mm in bilateral stance; Fig.~\ref{fig:loadprofiles}D). This asymmetry is accompanied by frontal pelvic tilt (abduction $\theta_y = 1.582^\circ$) and vertical shear ($\Delta\text{CC} = +0.394$~mm), as the unsupported contralateral hemi-pelvis drops while the loaded hemi-pelvis shears superiorly. This marked mechanical disparity is consistent with clinical observations that single-limb weight-bearing tasks (such as single-leg stance, stair ascent, and running) frequently provoke symptoms in patients with sacroiliac joint dysfunction and pelvic girdle instability \citep{Goode2008,Hammer2019,Keizer2019}.

\subsection*{Directional kinematic contrasts and mechanical proxies of parturition}
A key observation of this study is that evaluating SIJ response solely by scalar displacement magnitude provides an incomplete picture of joint mechanics. While parturition-motivated load configurations (LAB1--LAB3) produce lower scalar translation magnitudes ($0.66$--$0.79$~mm) than bipedal standing ($1.30$~mm), decomposition into local anatomical components reveals directional re-routing (Fig.~\ref{fig:directional}). In contrast to standing, internal ring compression reduces sagittal nutation ($\theta_{\text{nut}}$ decreases to $0.41^\circ$--$0.43^\circ$) and diminishes craniocaudal shear ($\Delta\text{CC} = -0.85$ to $-0.90$~mm). Instead, joint displacements are directed primarily into transverse outward distraction ($\Delta\text{ML} = +0.09$ to $+0.15$~mm) and sagittal gliding ($\Delta\text{AP} = +0.16$~mm).

This directional shift provides a mechanical perspective relevant to evolutionary discussions of the obstetrical dilemma \citep{Washburn1960,Ruff2017,Pavlicev2020}. To accommodate the passage of a neonate, the female pelvic ring must allow transient transverse and anteroposterior cavity dilation \citep{BiomechPregnancy2022,Fuchs2013PLOS,Almeida2024Uterus}. If parturition forces primarily engaged vertical shear or excessive counternutation, articular cartilages and dorsal ligaments would experience severe localized strain. By channeling internal expansion forces into transverse compliance ($\Delta\text{ML}$) while reducing vertical shear, the passive architecture of the pelvic ring facilitates outward dilation of the cavity while maintaining joint congruency. We emphasize, however, that the static forces applied here serve as standardized mechanical proxies and do not capture the dynamic, contact-driven progression of the fetal head through the birth canal.

\subsection*{Sex-related differences and morphometric associations}
Multivariable regression models demonstrated a statistically significant female translational surplus across all labor configurations ($+0.19$ to $+0.23$~mm in fully adjusted models, $p \le 0.020$; Fig.~\ref{fig:sexeffects}). Importantly, our nested regression analysis showed that this surplus attenuates by approximately 38\% upon adjusting for total pelvic volume and the morphometric triad (AP diameter, biischiadic width, and subpubic arch angle). This indicates that a substantial component of the observed sex-related difference is explained by sexual dimorphism in pelvic size and anterior arch architecture.

Furthermore, compliance in LAB1 was positively associated with the subpubic arch angle (Spearman $r = 0.52, p = 2.2 \times 10^{-20}$; Fig.~\ref{fig:sexeffects}C). Anatomically, the subpubic arch forms the anterior arch of the lower pelvic ring \citep{Ruff2017,Betti2020,Weiss2020}. A broader subpubic angle ($85.8^\circ$ in females vs.\ $69.1^\circ$ in males; Table~\ref{tab:cohort}) increases the transverse span relative to the pubic symphysis, decreasing the hoop stiffness of the anterior ring. Consequently, internal transverse expansion forces produce greater bilateral displacement at the dorsal SIJs in female pelves. This architectural relationship provides a structural hypothesis for future investigation of peripartum pelvic girdle pain (PGP), where broader pelvic dimensions may coincide with greater passive compliance under mechanical loading \citep{Hammer2019,Keizer2019}.

\subsection*{Skeletal geometry dominates bone material heterogeneity}
A central methodological finding of this work is the variance-channel comparison across matched simulation cohorts (Fig.~\ref{fig:variance}). Between-subject dispersion in SIJ micromotion was governed predominantly by macroscopic skeletal geometry: the shape-only model accounted for $98.7\%$ to $102.1\%$ (median $100.5\%$) of full-model population variance, with paired subject-level diagnostics showing near-perfect agreement ($R^2 \ge 0.994$, Pearson $r \ge 0.997$, $\text{RMSE} \le 0.031$~mm). In contrast, the material-only model accounted for less than $0.32\%$ (median $0.033\%$) of population variance.

This disparity has a solid basis in structural mechanics. The human pelvis functions as an indeterminate closed arch ring. The flexural and torsional compliances of an arch ring scale with the cube of cortical thickness and the geometry of the cross-section ($I \propto b\cdot h^3$), whereas material modulus $E$ enters linearly in the denominator. Across the adult population, inter-individual variations in pelvic dimensions, cortical curvatures, and articular orientations vary sufficiently to dominate load routing, rendering the influence of local bone mineral density heterogeneity secondary under standardized external loading.

Importantly, this comparison specifically isolates *bone mineral density (modulus) heterogeneity* derived from clinical CT scans, while keeping cartilage and ligament stiffness nominal. These results suggest that for macroscopic joint kinematic predictions, accurate subject-specific anatomical segmentation is paramount, whereas simplified bone density distributions introduce minimal bias into kinematic endpoints.

\subsection*{Secondary allometric size scaling}
Secondary log--log regressions indicated that allometric size scaling is load-dependent (Table~\ref{tab:allometry}; Fig.~\ref{fig:allometry}). Statistically significant negative allometric volume scaling was observed during single-leg stance (SP1leg: $\beta_{\text{male}} = -0.57, \beta_{\text{female}} = -0.43, q \le 0.004$ for translation), indicating that larger pelvic volume provides greater structural stiffness against asymmetric vertical shear. Under two-leg stance and parturition configurations, volume scaling exponents were generally attenuated, although female translation in LAB2 retained a statistically significant negative slope ($\beta = -0.614, q = 0.0360$).

Importantly, sex-by-size interaction terms were not statistically significant across any loading regime ($q \ge 0.958$), indicating that male and female pelves follow consistent scaling trajectories. Sex-related differences in SIJ compliance thus reflect an intercept shift related to dimorphic arch geometry rather than distinct scaling exponents.

\subsection*{Biomechanical and clinical considerations}
These findings provide mechanical context for orthopedic and pelvic health considerations:
\begin{enumerate}
\item \textbf{Sacroiliac Joint Arthrodesis:} Minimally invasive SIJ fusion is utilized for chronic, recalcitrant joint pain. Our data indicate that the greatest kinematic demand occurs during unilateral stance, marked by a substantial increase in bilateral shear and frontal tilt. These findings suggest that surgical fixation strategies should prioritize resistance to asymmetric craniocaudal shear and torsional tilt rather than symmetric axial compression \citep{Kiapour2019,Anderson2020}.
\item \textbf{Pelvic Girdle Compliance:} The positive correlation between subpubic arch angle and transverse compliance ($r = 0.52$) highlights the role of anterior ring geometry in posterior joint mechanics. While external pelvic compression belts are clinically applied to enhance force closure, further biomechanical studies are needed to determine whether individuals with broader arch dimensions experience differential benefit from external bracing \citep{Vleeming2012,Hammer2013}.
\end{enumerate}

\subsection*{Strengths and limitations}
The strengths of this study include the evaluation of a standardized 278-subject finite-element cohort across an adult lifespan, complete morphometric profiling, matched simulation variants isolating shape from bone material heterogeneity, and robust multivariable modeling with multiplicity control.

Several limitations must be acknowledged. First, simulations were formulated under static small-strain continuum elastostatics with linear-elastic bone, cartilage, and symphysis representations; dynamic inertial effects during gait and viscoelastic soft-tissue relaxation were not modeled. Second, the parturition load cases represent standardized static mechanical proxies rather than dynamic simulations of fetal head engagement and soft-tissue contact \citep{ChenGrimm2021,Almeida2024Uterus,BiomechPregnancy2022}. Third, the models did not incorporate pregnancy-induced hormonal laxity (e.g., relaxin-mediated ligament softening), which would be expected to augment compliance in vivo. Consequently, the observed female surplus reflects baseline skeletal and passive ligamentous architecture. Incorporating subject-specific soft-tissue properties and dynamic contact formulations represents an important area for future research.
